\begin{terminologyList}
    \term[Алфавит $\Sigma_\sigma$]{конечное непустое множество символов фиксированного размера $\sigma \geq 2$; в работе используется $\Sigma_\sigma = \{0, 1, \ldots, \sigma-1\}$}
    \term[Валидная маска]{целое число $x \in [0, 2^{N \cdot \mathrm{BITS}})$, в котором каждый из $N$ блоков по $\mathrm{BITS} = \lceil \log_2 \sigma \rceil$ бит кодирует символ из $\Sigma_\sigma$ (т.\,е. лежит в $[0, \sigma)$); множество валидных масок имеет размер $\sigma^N$}
    \term[Динамическое программирование (ДП, Dynamic Programming, DP)]{метод решения задач путём разложения их на подзадачи и запоминания решений подзадач для повторного использования}
    \term[Константа Чватала--Санкоффа $\gamma_\sigma$]{предел отношения $\mathbb{E}[L_n]/n$ при $n \to \infty$, где $L_n$~--- длина наибольшей общей подпоследовательности двух независимых равномерно случайных строк длины $n$ над алфавитом размера $\sigma$}
    \term[Лемма Фекете]{утверждение о том, что для любой супераддитивной последовательности $\{g_n\}$ (т.\,е. $g_{n+m} \geq g_n + g_m$) существует предел $\lim_{n \to \infty} g_n / n = \sup_n g_n / n \in (-\infty, +\infty]$}
    \term[Маска (состояния окна)]{целочисленное упакованное представление $N$-символьного окна одной из строк: $a = \sum_{p=0}^{N-1} a_p \cdot 2^{p \cdot \mathrm{BITS}}$, где $a_p \in \Sigma_\sigma$ --- $p$-й символ окна (нумерация от младшей позиции)}
    \term[Метод Monte-Carlo]{метод приближённого вычисления математических величин с помощью многократной случайной выборки и статистического усреднения; в работе используется для эмпирической оценки $\gamma_\sigma$ через классический ДП-LCS}
    \term[Наибольшая общая подпоследовательность (LCS, Longest Common Subsequence)]{максимальная по длине последовательность символов, которая является подпоследовательностью каждой из двух (или более) данных строк; длина LCS пары $(A, B)$ обозначается $\mathrm{lcs}(A, B)$}
    \term[Окно (look-ahead window)]{в данной работе~--- упорядоченный набор из $N$ ближайших ещё не использованных символов каждой из двух строк, видимый стратегии в текущий момент; в реализации кодируется маской (см.\,\S2.2)}
    \term[Подпоследовательность]{строка, полученная из данной вычёркиванием некоторого (возможно, пустого) подмножества символов с сохранением порядка оставшихся}
    \term[Сдвиг ($\mathrm{shift}$, $\mathrm{sh}$, $\mathrm{SH}$)]{разность $j - i$ позиций, на которых в данный момент стоят указатели по второй и первой строкам; в работе ограничен значениями из отрезка $[0,\,\mathrm{MAX\_SH}]$}
    \term[Состояние ДП]{в данной работе~--- тройка $(\mathrm{sh}, a, b)$, где $\mathrm{sh} \in [0, \mathrm{MAX\_SH}]$~--- текущий сдвиг, $a$ и $b$~--- маски следующих $N$ символов первой и второй строк соответственно; множество допустимых состояний обозначается $\mathcal{S}$}
    \term[Стратегия с ограниченным обзором (стратегия $\pi$)]{детерминированный или случайный алгоритм построения общей подпоследовательности двух строк, выбирающий на каждом шаге одно из трёх действий --- сдвиг по первой строке (L), сдвиг по второй строке (R) или засчитывание совпадения младших символов окон --- как функцию текущего состояния $(\mathrm{sh}, a, b)$ (см.\,\S2.1)}
    \term[Супераддитивность]{свойство последовательности $\{g_n\}$, состоящее в выполнении неравенства $g_{n+m} \geq g_n + g_m$ для всех $n, m \geq 0$; обеспечивает применимость леммы Фекете}
    \term[$f_K(\mathrm{sh}, a, b)$]{максимальное ожидаемое число совпадений за $K$ шагов из состояния $(\mathrm{sh}, a, b)$ при оптимальной стратегии; центральная ДП-функция работы, определяется в~\S2.3}
    \term[AVX2 (Advanced Vector Extensions 2)]{расширение набора инструкций x86, поддерживающее 256-битные векторные операции с целыми числами и числами с плавающей точкой}
\end{terminologyList}
